\documentclass[11pt,a4paper]{article}
\usepackage[margin=2.5cm]{geometry}
\usepackage[T1]{fontenc}
\usepackage{booktabs}
\usepackage{xcolor}
\usepackage{listings}
\usepackage[hidelinks]{hyperref}

\emergencystretch=2em
\definecolor{codebg}{gray}{0.95}
\lstset{language=C++, basicstyle=\ttfamily\small, keywordstyle=\color{blue!70!black},
  commentstyle=\color{green!50!black}\itshape, backgroundcolor=\color{codebg},
  frame=single, rulecolor=\color{black!20}, breaklines=true, showstringspaces=false,
  tabsize=2, columns=flexible}

\title{ATmega32A Driver Course \\ \Large Module 7 --- SPI (Serial Peripheral Interface)}
\author{Ali Sahafi}
\date{}

\begin{document}
\maketitle

\section{Theory}

SPI is a fast, synchronous bus for talking to peripherals such as displays,
SD cards, sensors and shift registers. One \textbf{master} generates the
clock; one or more \textbf{slaves} respond. Four wires:

\begin{center}
\begin{tabular}{llll}
\toprule
\textbf{Signal} & \textbf{Pin} & \textbf{Direction (master)} & \textbf{Meaning} \\
\midrule
MOSI & PB5 & output & Master Out, Slave In \\
MISO & PB6 & input  & Master In, Slave Out \\
SCK  & PB7 & output & Clock (generated by master) \\
SS   & PB4 & output & Slave Select (active low) \\
\bottomrule
\end{tabular}
\end{center}

SPI is built around a \textbf{shift register}: with every clock pulse one bit
goes out on MOSI \emph{and} one bit comes in on MISO. Sending and receiving
are the same operation --- \texttt{transfer()} always returns the byte the
slave shifted back. To ``only read'', you send a dummy byte.

\paragraph{Modes.} Clock polarity (CPOL) and phase (CPHA) give four modes
(\texttt{SPI\_MODE\_0} \dots\ \texttt{SPI\_MODE\_3}); the datasheet of the
slave device tells you which one to use. Mode 0 is by far the most common.

\section{Registers}

\begin{center}
\small
\begin{tabular}{lll}
\toprule
\textbf{Register} & \textbf{Role} & \textbf{Bits used by the driver} \\
\midrule
\texttt{SPCR} & Control & \texttt{SPE}, \texttt{MSTR}, \texttt{CPOL}, \texttt{CPHA}, \texttt{DORD}, \texttt{SPR1:0}, \texttt{SPIE} \\
\texttt{SPSR} & Status  & \texttt{SPIF} (transfer complete), \texttt{SPI2X} (double speed) \\
\texttt{SPDR} & Data    & write to start a transfer, read to get the answer \\
\bottomrule
\end{tabular}
\end{center}

\section{API reference}

\begin{center}
\footnotesize
\begin{tabular}{ll}
\toprule
\textbf{Method} & \textbf{Description} \\
\midrule
\texttt{SPI.initMaster(mode, prescaler, order, irq)} & Configure as SPI master \\
\texttt{SPI.initSlave(mode, order, irq)} & Configure as SPI slave \\
\texttt{SPI.transfer(data)} & Synchronous transfer (returns received byte) \\
\texttt{SPI.attachInterrupt(callback)} & Transfer-complete callback: \texttt{void cb(uint8\_t)} \\
\bottomrule
\end{tabular}
\end{center}

\begin{itemize}
  \item \textbf{Mode:} \texttt{SPI\_MODE\_0/1/2/3}
  \item \textbf{Prescaler:} \texttt{SPI\_PRESCALER\_4/16/64/128}
  \item \textbf{Double-speed prescaler:}
        \texttt{SPI\_PRESCALER\_2\_2X / 8\_2X / 32\_2X / 64\_2X}
  \item \textbf{Data order:} \texttt{SPI\_MSB\_FIRST}, \texttt{SPI\_LSB\_FIRST}
\end{itemize}

\paragraph{Pins are configured automatically.} \texttt{initMaster()} and
\texttt{initSlave()} set the DDR bits for PB4--PB7 correctly; you do not
touch \texttt{DDRB} yourself. As master, keep PB4 (SS) an output --- if it
floats and reads LOW, the hardware silently drops out of master mode.

\section{Example: loopback test}

The simplest possible SPI experiment needs no slave at all: connect
\textbf{MOSI to MISO with one jumper wire} and every transmitted byte comes
straight back.

\paragraph{Wiring.}
\begin{itemize}
  \item Jumper wire from \textbf{PB5 (MOSI)} to \textbf{PB6 (MISO)}.
  \item UART on PD0/PD1 at 9600 baud to watch the results (Module 6).
\end{itemize}

\paragraph{Build \& flash.} \texttt{make spi}

\lstinputlisting{example.cpp}

Remove the jumper and every read returns a constant --- there is nothing
driving MISO anymore.

\section{Exercises}

\begin{enumerate}
  \item Program a second board with \texttt{SPI.initSlave()} and an interrupt
        callback that stores the received byte; wire the two boards together
        and exchange a counter.
  \item Connect a 74HC595 shift register (SER $\leftarrow$ MOSI, SRCLK
        $\leftarrow$ SCK, RCLK $\leftarrow$ SS) and drive 8 LEDs over SPI ---
        three wires instead of eight.
  \item Change the data order to \texttt{SPI\_LSB\_FIRST} on the master only.
        What arrives at the slave when you send \texttt{0x01}? Verify by
        experiment.
\end{enumerate}

\end{document}
